\chapter*{BẢNG PHÂN CÔNG NHIỆM VỤ NHÓM THỰC HIỆN}
\addcontentsline{toc}{chapter}{BẢNG PHÂN CÔNG NHIỆM VỤ NHÓM THỰC HIỆN}
\markboth{PHÂN CÔNG NHIỆM VỤ}{PHÂN CÔNG NHIỆM VỤ}

Hai thành viên cùng thực hiện các nội dung nghiên cứu, thiết kế, triển khai và biên soạn báo cáo. Phần thực nghiệm được phân công theo thiết bị: \reportstudentone{} phụ trách cấu hình i5, \reportstudenttwo{} phụ trách cấu hình i7. Hai thành viên phối hợp rà soát kết quả và chịu trách nhiệm chung về nội dung báo cáo.

\begin{table}[htbp]
\centering
\caption*{Phân công nhiệm vụ nhóm thực hiện đề tài}
\small
\renewcommand{\arraystretch}{1.28}
\setlength{\tabcolsep}{6pt}
\begin{tabularx}{\textwidth}{@{}>{\RaggedRight\arraybackslash}p{5.0cm} >{\RaggedRight\arraybackslash}p{3.8cm} >{\RaggedRight\arraybackslash}X@{}}
\toprule
\rowcolor{tableheader}
\textbf{Nhóm nhiệm vụ} & \textbf{Thành viên phụ trách} & \textbf{Kết quả thực hiện} \\
\midrule
Nghiên cứu Transformer, MoE, Router và phân tầng bộ nhớ & \reportstudentone\newline\reportstudenttwo & Cơ sở lý thuyết và mô tả bài toán \\
\addlinespace[5pt]
Phân tích, thiết kế và triển khai nguyên mẫu suy luận & \reportstudentone\newline\reportstudenttwo & Kiến trúc và các thành phần hệ thống \\
\addlinespace[5pt]
Khảo sát cache, offloading và prefetching & \reportstudentone\newline\reportstudenttwo & Thuật toán và cơ chế quản lý expert \\
\addlinespace[5pt]
Thực nghiệm và đánh giá vận hành mô hình trên cấu hình i5 & \reportstudentone & Hiệu năng, bộ nhớ và nhận xét cấu hình i5 \\
\addlinespace[5pt]
Thực nghiệm và đánh giá vận hành mô hình trên cấu hình i7 & \reportstudenttwo & Hiệu năng, bộ nhớ và nhận xét cấu hình i7 \\
\addlinespace[5pt]
Tổng hợp kết quả, biên soạn và rà soát báo cáo & \reportstudentone\newline\reportstudenttwo & Báo cáo định kỳ lần 1 \\
\bottomrule
\end{tabularx}
\end{table}
